Quantization is one of the most successful methods for optimizing neural networks for efficient inference and on-device execution while maintaining a high accuracy. By compressing the weights and activations from the regular 32-bit floating-point format to more efficient low bit fixed-point representations, such as INT8, we can reduce power consumption and accelerate inference when deploying neural networks on edge devices \citep{horowitz}.

% reduce the memory footprint  and accelerate inference, thus, lowering the power consumption and silicon area when deploying neural networks on edge devices \citep{horowitz}. 

Despite its clear power and latency benefits, quantization comes at the cost of added noise due to the reduced precision. However, researchers in recent years have shown that neural networks are robust to this noise and can be quantized to 8-bits with a minimal drop in accuracy using \textit{post-training quantization} techniques (PTQ) \citep{dfq, bannerposttraining, nahshan2020lapq}. PTQ can be very efficient and, generally, only requires access to a small calibration dataset, but suffers when applied to low-bit quantization ($\leq$ 4-bits) of neural networks. Meanwhile, \textit{quantization-aware training} (QAT) has become the de-facto standard method for achieving low-bit quantization while maintaining near full-precision accuracy \citep{krishnamoorthi,lsq, whitepaper}. By simulating the quantization operation during training or fine-tuning, the network can adapt to the quantization noise and reach better solutions than with PTQ.  
 
% Exploiting the inherent resilience of neural networks to random perturbations, researchers in recent years have shown that neural networks can be quantized to 8-bits with minimal drop in accuracy using \textit{post-training quantization} techniques (PTQ) \citep{dfq, bannerposttraining, nahshan2020lapq}. 


In this paper, we focus on the \textit{oscillations} of quantized weights that occur during quantization-aware training. This is a little-known and under-investigated phenomenon in the optimization of quantized neural networks, with significant consequences for the network during and after training. When using the popular \textit{straight-through estimator} (STE) \citep{bengio2013estimating} for QAT, weights seemingly randomly oscillate between adjacent quantization levels leading to detrimental noise during the optimization process. Equipped with this insight, we investigate recent advances in QAT that claim improved performance and assess their effectiveness in addressing this oscillatory behavior.

An adverse symptom of weight oscillations is that they can corrupt the estimated inference statistics of the batch-normalization layer collected during training, leading to poor validation accuracy. We find that this effect is particularly pronounced in low-bit quantization of efficient networks with depth-wise separable layers, such as MobileNets or EfficientNets, but can be addressed effectively by re-estimating the batch-normalization statistics after training. %\citep{louizos2018relaxed, peters2018probabilistic}.

While batch-normalization re-estimation overcomes one significant symptom of the oscillations, it does not tackle its root cause. To this end, we propose two novel algorithms that are effective at reducing oscillations: \textit{oscillation dampening} and  \textit{iterative weight freezing}. By addressing oscillations at their source, our methods improve accuracy beyond the level of batch-normalization re-estimation. We show that both methods achieve state-of-the-art results for 4 and 3-bit quantization of efficient networks, such as \mnv{2}, \mnv{3}, and \efnet{lite} on ImageNet.

